\section{État de l'art}

\subsection{Positionnement de l'état de l'art}

Le problème posé par ce mémoire n'est pas seulement celui de l'utilisation d'un chatbot
ou d'un générateur de texte. Il concerne la conception d'une plateforme de formation
autonome alimentée par l'intelligence artificielle. Deux dimensions doivent donc être
étudiées conjointement. La première concerne l'assistant pédagogique : comment permettre
à un modèle de langage de répondre aux questions des apprenants à partir des contenus
réels de la formation ? La seconde concerne la production de contenus : comment générer
des cours longs, structurés, oralisables et contrôlables sans se limiter à une simple sortie
textuelle produite par un prompt ?

L'état de l'art est volontairement centré sur les techniques plutôt que sur un catalogue
d'outils commerciaux. Les plateformes d'apprentissage en ligne permettent déjà de
diffuser des cours, de suivre des élèves ou de proposer des évaluations, mais elles ne
répondent pas directement à la question de recherche de ce mémoire : rendre une chaîne
de formation suffisamment autonome tout en conservant une qualité pédagogique, une
traçabilité et une capacité d'évaluation. Les approches étudiées ici sont donc celles qui
permettent d'adapter ou d'encadrer les modèles de langage dans un contexte métier.

La progression adoptée est la suivante. On étudie d'abord les approches permettant de
donner un contexte à un modèle de langage : prompt engineering, fine-tuning et RAG. On
analyse ensuite les variantes du RAG, notamment la recherche vectorielle, la recherche
hybride, le reclassement sémantique, les métadonnées et l'évaluation. Enfin, on étudie
les méthodes de génération contrôlée de contenus longs : décomposition de tâche,
génération structurée, schémas intermédiaires, auto-évaluation et validation humaine. Ce
cheminement prépare l'étude originale, sans encore détailler la solution développée.

\subsection{Modèles de langage et adaptation au contexte}

Les grands modèles de langage ont montré qu'ils peuvent réaliser de nombreuses tâches
à partir d'instructions en langage naturel et de quelques exemples fournis dans le
prompt. Les travaux sur le few-shot learning ont mis en évidence cette capacité à adapter
le comportement du modèle sans modifier ses poids \cite{brown2020language}. Pour une
plateforme de formation, cette propriété est importante : un modèle généraliste peut être
orienté vers un rôle d'assistant pédagogique, de rédacteur de cours ou de correcteur de
contenu simplement à partir d'instructions.

Cette flexibilité explique l'intérêt du prompt engineering. Dans sa forme la plus simple,
il consiste à fournir au modèle une consigne, un rôle, des contraintes de style, un format
de sortie et parfois des exemples. Dans le cas d'un chatbot pédagogique, le prompt peut
indiquer que le modèle doit répondre clairement, ne pas inventer d'information et rester
dans le périmètre du cours. Dans le cas d'une génération de contenu, il peut préciser le
niveau de l'apprenant, le ton attendu, la structure du cours ou les règles de conformité.

Le prompt engineering présente trois avantages. Il est rapide à mettre en place, peu
coûteux à modifier et compatible avec des modèles propriétaires accessibles par API. Il
est donc très utile en phase de prototypage. Cependant, il a aussi des limites fortes. Le
contexte disponible dans le prompt reste limité ; les consignes peuvent entrer en
concurrence ; et le modèle peut suivre certaines instructions tout en en négligeant
d'autres. Lorsque le volume documentaire augmente ou lorsque les règles métier deviennent
nombreuses, le prompt devient difficile à maintenir.

Surtout, le prompt ne constitue pas une mémoire documentaire fiable. Il permet de
fournir ponctuellement des informations, mais il ne garantit pas que le modèle réponde
à partir d'une source précise. Dans une formation, cette limite est critique : l'assistant
doit répondre à partir des contenus réellement enseignés, pas à partir d'une connaissance
générale ou approximative. De même, un générateur de cours doit respecter un programme
et ne pas produire une progression plausible mais non conforme.

\subsection{Fine-tuning et spécialisation du modèle}

Une autre approche consiste à adapter le modèle par fine-tuning. Au lieu de fournir toutes
les règles dans le prompt, on entraîne le modèle sur des exemples représentatifs de la
tâche attendue. Cette approche peut améliorer le style, le format de sortie ou la
répétabilité d'un comportement. Elle est particulièrement utile lorsque l'on souhaite
spécialiser un modèle sur un type de réponse ou un domaine récurrent.

Le fine-tuning classique peut cependant être coûteux, car il nécessite de modifier un
grand nombre de paramètres. Des méthodes plus légères, comme LoRA, cherchent à
réduire ce coût en gelant les poids du modèle et en ajoutant des matrices de bas rang
entraînables \cite{hu2021lora}. Ces approches rendent la spécialisation plus accessible,
mais elles ne suppriment pas les contraintes fondamentales du fine-tuning : il faut
constituer un jeu de données de qualité, entraîner ou adapter un modèle, évaluer le
résultat et maintenir cette adaptation dans le temps.

Dans le contexte d'un assistant pédagogique, le fine-tuning n'est pas toujours la meilleure
réponse au problème de connaissance. Il permet d'apprendre un style ou une procédure,
mais il est moins adapté pour intégrer des documents qui changent souvent. Si un cours,
un référentiel ou un support est modifié, il n'est pas souhaitable de réentraîner le modèle
à chaque mise à jour. En outre, le fine-tuning ne garantit pas que la réponse produite soit
fidèle à un document précis. Un modèle fine-tuné peut continuer à halluciner ou à
mélanger des informations apprises avec des informations plausibles.

Pour une plateforme de formation, le fine-tuning est donc pertinent pour stabiliser un
comportement, mais insuffisant comme mécanisme principal de contextualisation
documentaire. Il peut apprendre au modèle à répondre comme un assistant pédagogique,
mais il ne suffit pas à lui donner une mémoire fiable, actualisable et traçable des cours.

\subsection{RAG : ancrer la génération dans des documents}

Le Retrieval-Augmented Generation, ou RAG, répond à cette limite en séparant deux
fonctions : la récupération d'information et la génération de réponse. Le principe est de
ne pas modifier directement le modèle, mais de rechercher au moment de la question des
passages pertinents dans une base documentaire, puis de fournir ces passages au modèle
comme contexte. Le modèle génère alors une réponse à partir des documents récupérés.
Cette approche a été formalisée dans les travaux de Lewis et al. sur les tâches de NLP
intensives en connaissances \cite{lewis2020retrieval}.

Le RAG présente plusieurs avantages pour un assistant pédagogique. D'abord, les contenus
peuvent être mis à jour sans réentraîner le modèle. Il suffit de réindexer les documents.
Ensuite, la réponse peut être rattachée à des passages sources, ce qui améliore la
traçabilité. Enfin, le modèle est moins dépendant de sa mémoire paramétrique : il peut
répondre à partir du cours réellement fourni, même si ce contenu est spécifique à une
formation ou à une promotion.

Cette architecture est particulièrement adaptée au contexte de ce mémoire. Les élèves ne
doivent pas recevoir une réponse générique sur un métier ou une notion ; ils doivent
recevoir une réponse cohérente avec le cours diffusé, le référentiel utilisé et les supports
mis à disposition. Le RAG permet donc de transformer un modèle généraliste en assistant
contextualisé.

Cependant, le RAG ne supprime pas tous les risques. Il déplace une partie du problème
vers la qualité de la récupération. Si les passages récupérés sont incomplets, trop larges,
hors sujet ou mal classés, le modèle peut produire une réponse faible. De plus, même
avec un bon contexte, le modèle peut reformuler de manière imprécise ou ajouter une
information non supportée par les documents. Les travaux sur les hallucinations montrent
que les modèles de langage peuvent générer des informations plausibles mais non fondées,
ce qui impose des mécanismes d'évaluation et de contrôle \cite{ji2023hallucination}.

\subsection{Techniques de recherche dans les architectures RAG}

Un système RAG n'est pas une technique unique. Il regroupe plusieurs choix d'architecture
qui influencent directement la qualité des réponses. Le premier choix concerne le
découpage documentaire. Les documents sont rarement injectés entiers dans le modèle :
ils sont découpés en segments, ou chunks, qui seront indexés puis récupérés. Des chunks
trop petits peuvent perdre le contexte nécessaire à la compréhension ; des chunks trop
grands peuvent diluer l'information pertinente. Le choix de la taille, du chevauchement
et des métadonnées est donc déterminant.

Le deuxième choix concerne la représentation des documents. Une recherche par mots-clés
fonctionne bien lorsque la question reprend le vocabulaire exact du document, mais elle
peut échouer lorsque l'élève reformule avec d'autres termes. La recherche vectorielle
répond à cette limite en représentant les questions et les passages dans un espace
sémantique. Elle permet de retrouver des passages proches par le sens, même si les mots
ne sont pas identiques.

La recherche vectorielle a elle aussi des limites. Elle peut retrouver des passages
sémantiquement proches mais pas exactement pertinents, ou négliger des mots clés
importants comme une référence, un code, une notion réglementaire ou un terme métier.
C'est pourquoi de nombreuses architectures combinent recherche lexicale et recherche
vectorielle. On parle alors de recherche hybride. Les travaux récents sur le RAG décrivent
cette évolution d'un RAG naïf vers des architectures plus avancées ou modulaires
\cite{gao2023retrieval}.

Le troisième choix concerne le reclassement des résultats. Après une première recherche,
un système peut reranker les passages candidats pour privilégier ceux qui correspondent
le mieux à l'intention de la question. Ce reclassement peut s'appuyer sur des modèles de
langage ou des systèmes spécialisés. Dans un contexte pédagogique, il est utile car les
questions des élèves sont souvent courtes, imprécises ou formulées avec des termes non
experts.

Le quatrième choix concerne les métadonnées. Un même corpus peut contenir plusieurs
formations, plusieurs jours, plusieurs modules ou plusieurs types de documents. Sans
filtrage, l'assistant risque de récupérer un passage provenant d'un mauvais cours. Les
métadonnées permettent de restreindre la recherche à un périmètre pertinent : formation,
jour, module, type de support ou version du document.

\subsection{Azure AI Search comme approche industrielle du RAG}

Plusieurs solutions peuvent être utilisées pour construire un RAG : base vectorielle locale,
index spécialisé, moteur de recherche hybride ou service cloud managé. Dans un prototype,
une bibliothèque locale comme FAISS peut suffire pour tester la recherche vectorielle.
Dans une plateforme destinée à être utilisée par de vrais élèves, d'autres critères
deviennent importants : indexation automatique, recherche hybride, filtrage, supervision,
sécurité, intégration avec le stockage documentaire et capacité à maintenir plusieurs
corpus.

Azure AI Search s'inscrit dans cette logique industrielle. La documentation Microsoft
présente Azure AI Search comme un service d'indexation et de recherche pouvant être
utilisé dans des architectures RAG avec des données textuelles et vectorielles
\cite{microsoft2026foundryrag}. Le service permet de combiner recherche par mots-clés
et recherche vectorielle dans une requête hybride ; les résultats sont fusionnés avant
d'être exploités par l'application \cite{microsoft2026hybridsearch}. Il propose également
un semantic ranker, qui reclasse les résultats textuels ou hybrides à partir de modèles de
compréhension du langage \cite{microsoft2026semanticranker}.

Dans le contexte d'un chatbot pédagogique, ces fonctionnalités répondent à plusieurs
besoins. La recherche hybride permet de traiter à la fois les formulations approximatives
des élèves et les termes précis présents dans les supports. Le semantic ranking permet
d'améliorer la pertinence des passages présentés au modèle. Les index et métadonnées
permettent de séparer les formations ou les plateformes. Enfin, l'intégration avec Azure
OpenAI simplifie la construction d'une chaîne de question-réponse où les documents
indexés sont fournis comme source au modèle.

Cette approche a toutefois des limites. Un service managé réduit la complexité
d'infrastructure, mais il ne dispense pas de concevoir correctement le corpus. La qualité
du RAG dépend encore du découpage, des métadonnées, des documents indexés, de la
formulation des requêtes et du prompt final. Le semantic ranker reclasse des résultats
déjà récupérés ; il ne peut pas compenser entièrement une mauvaise indexation ou un
corpus incomplet. Le choix d'Azure AI Search doit donc être compris comme un choix
d'architecture robuste, pas comme une garantie automatique de qualité.

\subsection{Évaluation des systèmes RAG}

L'évaluation d'un RAG doit distinguer la qualité de la récupération et la qualité de la
réponse. Une réponse peut être bien formulée mais fondée sur de mauvais passages ; à
l'inverse, les bons passages peuvent être récupérés mais mal exploités par le modèle.
Plusieurs cadres d'évaluation récents proposent donc de mesurer séparément la pertinence
du contexte, la fidélité de la réponse aux sources et la pertinence de la réponse finale.
C'est le cas de RAGAS \cite{es2023ragas} et d'ARES \cite{saadfalcon2023ares}.

Dans un contexte pédagogique, les métriques utiles peuvent être regroupées en quatre
familles. La première concerne la récupération : précision@k, rappel@k, rang moyen du
bon passage, taux de questions pour lesquelles au moins une source pertinente est
retrouvée. La deuxième concerne la fidélité : la réponse est-elle supportée par les passages
fournis ? Ajoute-t-elle une information non présente dans les documents ? La troisième
concerne l'utilité pédagogique : la réponse est-elle claire, adaptée au niveau de l'élève et
suffisamment concise ? La quatrième concerne l'expérience réelle : les élèves reposent-ils
la même question, utilisent-ils l'assistant et jugent-ils les réponses utiles ?

La citation des sources est également un enjeu. Les travaux sur la vérifiabilité dans les
moteurs génératifs montrent que les systèmes doivent citer de manière complète et
précise les éléments qui soutiennent leurs réponses \cite{liu2023evaluating}. Pour une
formation, cette exigence est moins formelle que dans une recherche documentaire
académique, mais elle reste importante : un responsable pédagogique doit pouvoir
vérifier d'où vient une réponse.

Ces considérations conduisent à une conclusion claire : un RAG ne doit pas être évalué
uniquement par impression subjective. Il faut construire un jeu de questions représentatif,
identifier les passages attendus, comparer plusieurs configurations de recherche et
mesurer la fidélité des réponses. Sans cette évaluation, il est impossible de justifier
scientifiquement le choix d'une architecture RAG plutôt qu'une autre.

\subsection{Génération de contenus longs par prompt unique}

La deuxième dimension du mémoire concerne la génération automatique des contenus de
formation. La première approche possible consiste à demander directement à un modèle
de langage de produire un cours complet à partir d'un programme, d'un référentiel et de
quelques consignes. Cette approche est séduisante parce qu'elle exploite la capacité des
modèles à rédiger des textes cohérents et à suivre un format général.

Dans un contexte simple, la génération directe peut donner des résultats acceptables. Elle
est utile pour produire un premier brouillon, un résumé, un plan de cours ou une
explication courte. Elle montre rapidement que l'IA peut réduire le temps de rédaction.
Cependant, elle devient fragile lorsque la sortie attendue est longue, structurée,
oralisable et destinée à être utilisée sans réécriture complète.

Plusieurs problèmes apparaissent alors. Le modèle peut répéter des idées, perdre le fil de
la progression, produire des transitions artificielles, dépasser la longueur attendue ou
oublier certaines contraintes. Il peut aussi mélanger le langage interne du système avec
le langage destiné à l'apprenant, par exemple en verbalisant des contraintes techniques
qui ne devraient pas apparaître dans le cours. Enfin, si le résultat final est mauvais, il est
difficile de savoir quelle partie du prompt ou de la génération est responsable.

Ces limites ne signifient pas que le prompt engineering est inutile. Elles montrent plutôt
que le prompt unique atteint rapidement ses limites lorsqu'il doit porter toutes les
responsabilités : structure pédagogique, style oral, durée, conformité, exemples,
supports, exercices et qualité. Un prompt trop chargé devient un artefact fragile, difficile
à faire évoluer et difficile à évaluer.

\subsection{Décomposition de tâches et génération structurée}

Pour dépasser ces limites, de nombreux travaux sur les modèles de langage insistent sur
la décomposition des tâches complexes. Le chain-of-thought a montré que faire produire
des étapes intermédiaires peut améliorer certaines tâches de raisonnement
\cite{wei2022chain}. Le least-to-most prompting va plus loin en demandant au modèle de
décomposer un problème en sous-problèmes plus simples avant de les résoudre
\cite{zhou2022least}. Les approches de type Tree of Thoughts proposent d'explorer
plusieurs pistes intermédiaires avant de choisir une solution \cite{yao2023tree}.

Ces travaux portent principalement sur le raisonnement, mais leur logique est applicable
à la génération de contenus pédagogiques longs. Un cours complet est une tâche complexe
qui peut être décomposée : analyser le programme, construire un plan, détailler les
sections, rédiger les transitions, adapter le texte à l'oral, générer des supports, produire
des questions, puis vérifier la conformité. La qualité peut alors être contrôlée à chaque
étape plutôt qu'uniquement à la fin.

La génération structurée repose sur la même idée. Au lieu de produire directement un
texte final, le modèle produit d'abord une représentation intermédiaire : plan, JSON,
liste d'objectifs, sections, contraintes de durée ou points clés. Cette représentation sert
de contrat entre les étapes. Elle permet d'auditer la structure avant de rédiger, de répartir
les budgets, de générer section par section et de rattacher les supports à des intentions
pédagogiques explicites.

Cette approche présente des avantages importants. Elle réduit la charge imposée au
modèle à chaque appel, rend les erreurs plus localisables et facilite l'intégration logicielle.
Elle introduit aussi une forme de traçabilité : on peut comparer le plan prévu, le texte
généré et les corrections appliquées. En revanche, elle augmente la complexité de
l'orchestration. Il faut gérer plusieurs appels au modèle, valider les formats, traiter les
erreurs JSON, conserver les artefacts et maîtriser le coût.

\subsection{Auto-évaluation, reviews automatiques et validation humaine}

Une autre famille d'approches consiste à utiliser les modèles de langage non seulement
pour générer, mais aussi pour critiquer ou améliorer leurs propres sorties. Des travaux
comme Reflexion explorent l'idée d'agents capables d'utiliser un retour verbal pour
améliorer leurs actions \cite{shinn2023reflexion}. Dans les applications de génération de
contenu, cette idée se traduit souvent par des passes de review : un modèle relit le texte,
identifie des problèmes, propose des corrections ou vérifie le respect de règles.

Les reviews automatiques sont intéressantes pour une plateforme de formation, car elles
permettent de contrôler plusieurs dimensions : cohérence avec le plan, ton pédagogique,
adaptation à l'oral, conformité éthique, absence de formulations sensibles, respect de la
durée ou qualité des QCM. Elles peuvent réduire la charge de validation humaine en
signalant les passages problématiques.

Elles ne doivent toutefois pas être considérées comme infaillibles. Un modèle peut ne pas
détecter une erreur, proposer une correction qui dégrade le contenu ou valider une
réponse insuffisamment sourcée. L'évaluation par LLM-as-a-judge peut être utile, mais
elle possède des biais et doit être encadrée par des critères explicites \cite{zheng2023judging}.
Dans un système rigoureux, les reviews automatiques doivent donc être combinées avec
des règles déterministes, des métriques et, lorsque le risque le justifie, une validation
humaine.

La validation humaine garde une place importante dans les systèmes pédagogiques. Elle
peut intervenir à plusieurs niveaux : validation du programme, relecture d'échantillons,
contrôle des sorties à risque, analyse des retours élèves et amélioration des règles. Le
but n'est pas nécessairement de relire chaque phrase produite, mais de déplacer
l'intervention humaine vers les points où elle apporte le plus de valeur.

\subsection{Évaluation des contenus pédagogiques générés}

L'évaluation des contenus générés est plus difficile que l'évaluation d'une tâche fermée.
Un cours n'a pas une seule bonne réponse. Il faut donc combiner plusieurs indicateurs.
Le premier est la fidélité au programme : les compétences prévues sont-elles couvertes ?
Les notions importantes sont-elles présentes ? Le cours évite-t-il les informations hors
périmètre ?

Le deuxième indicateur est la cohérence pédagogique. Un cours doit suivre une progression
compréhensible, éviter les répétitions inutiles, introduire les notions avant de les utiliser
et proposer des exemples adaptés. Cette dimension peut être évaluée par relecture humaine,
par grille critériée ou par comparaison entre le plan et le texte généré.

Le troisième indicateur est l'adaptation à l'oral. Un texte destiné à la synthèse vocale ne
doit pas être écrit comme un support PDF. Il doit contenir des phrases prononçables, des
transitions, un rythme d'écoute et des formulations naturelles. La durée est également
un critère : si le texte dépasse le créneau audio prévu, la formation devient difficile à
diffuser.

Le quatrième indicateur concerne les supports associés. Les slides doivent correspondre
aux moments pédagogiques importants, et les QCM doivent vérifier des notions réellement
enseignées. Un QCM généré indépendamment du cours peut donner une illusion
d'évaluation sans mesurer la compréhension du contenu diffusé.

Enfin, il faut mesurer l'impact opérationnel. Une pipeline de génération n'a d'intérêt que
si elle réduit effectivement le temps de production, les manipulations manuelles ou le
coût, tout en maintenant une qualité acceptable. Les métriques doivent donc combiner
qualité pédagogique et efficacité industrielle.

\subsection{Synthèse critique}

Les approches étudiées montrent que les modèles de langage ne suffisent pas, à eux
seuls, à construire une plateforme de formation autonome. Le prompt engineering est
utile pour orienter le comportement du modèle, mais il devient fragile lorsque les règles
et les documents se multiplient. Le fine-tuning permet d'adapter un style ou une tâche,
mais il n'est pas idéal pour injecter des connaissances documentaires souvent mises à
jour. Le RAG apporte une solution plus adaptée pour l'assistant pédagogique, car il
ancre les réponses dans les contenus de formation et permet une mise à jour des sources.

De la même manière, la génération directe d'un cours par prompt unique est insuffisante
pour produire des contenus longs, oralisables et contrôlables. Les approches de
décomposition, de génération structurée et de review automatique offrent une meilleure
base pour construire une pipeline fiable. Elles transforment la génération en processus,
avec des étapes, des artefacts et des points de contrôle.

Le point commun entre ces deux axes est la nécessité d'encadrer l'IA. Dans un assistant
RAG, l'encadrement passe par la récupération, les sources et l'évaluation de la fidélité.
Dans une pipeline de génération, il passe par le plan, les formats structurés, les reviews
et les métriques. L'originalité de l'approche étudiée dans la suite du mémoire devra donc
être évaluée à partir de cette double exigence : contextualiser les réponses du chatbot et
structurer la génération des contenus.

\subsection{Tableau récapitulatif des approches}

\begin{table}[h]
\centering
\small
\begin{tabular}{|p{3cm}|p{3.4cm}|p{4.2cm}|p{3.4cm}|}
\hline
\textbf{Approche} & \textbf{Apport principal} & \textbf{Limites} & \textbf{Intérêt pour le projet} \\
\hline
Prompt engineering &
Orientation rapide du modèle par consignes, rôles, exemples et contraintes. &
Contexte limité, maintenance difficile, faible traçabilité, risque d'oubli de consignes. &
Utile pour le style et les règles locales, insuffisant seul pour une plateforme autonome. \\
\hline
Fine-tuning &
Adaptation du comportement ou du style du modèle sur des exemples. &
Coût de constitution des données, maintenance lourde, pas idéal pour documents changeants. &
Pertinent pour stabiliser un comportement, moins adapté à la connaissance pédagogique actualisable. \\
\hline
RAG simple &
Recherche de passages pertinents puis génération ancrée dans les documents. &
Dépend du chunking, de la qualité du corpus et de la pertinence de la récupération. &
Base pertinente pour répondre aux élèves à partir des contenus de cours. \\
\hline
RAG avancé &
Recherche hybride, reranking, métadonnées, citations et évaluation. &
Architecture plus complexe, coûts supplémentaires, besoin de jeux de tests. &
Approche la plus adaptée pour un assistant pédagogique traçable et maintenable. \\
\hline
Azure AI Search &
Service managé combinant indexation, recherche vectorielle, recherche hybride et semantic ranker. &
Ne compense pas un mauvais corpus ; dépend de la configuration, des filtres et de l'évaluation. &
Choix cohérent pour industrialiser le RAG dans une plateforme utilisée par de vrais élèves. \\
\hline
Génération directe de cours &
Production rapide de brouillons ou de contenus courts. &
Dérive sur contenus longs, répétitions, contrôle faible de la durée et de la structure. &
Approche utile au prototypage, mais insuffisante comme chaîne de production. \\
\hline
Génération structurée &
Plan, sections, schémas et artefacts intermédiaires. &
Orchestration plus complexe, validation JSON et coût d'appels multiples. &
Approche adaptée à la production de cours longs, audios, slides et QCM. \\
\hline
Reviews automatiques &
Contrôle du style, de la conformité, de l'adhérence au plan et de la qualité. &
Risque de faux positifs ou faux négatifs ; nécessite métriques et supervision. &
Permet de réduire la relecture humaine tout en rendant la qualité plus observable. \\
\hline
Human-in-the-loop &
Validation humaine sur les points critiques. &
Réintroduit du coût humain et limite l'autonomie complète. &
Nécessaire pour superviser le système, surtout en production pédagogique. \\
\hline
\end{tabular}
\caption{Synthèse des approches existantes et de leur intérêt pour une plateforme de formation autonome}
\end{table}

\subsection{Comparaisons et évaluations à réaliser}

Pour rendre le mémoire scientifiquement rigoureux, plusieurs comparaisons doivent être
réalisées ou, à défaut, présentées comme des évaluations à consolider. La première
comparaison concerne le chatbot : prompt seul contre RAG. Il s'agit de construire un
jeu de questions d'élèves, de définir les passages sources attendus, puis de comparer la
fidélité, la pertinence et le taux d'hallucination des réponses.

La deuxième comparaison concerne les variantes de RAG. Il serait pertinent de comparer
une recherche vectorielle simple, une recherche hybride et une recherche hybride avec
reclassement sémantique lorsque la configuration le permet. Les métriques attendues
sont la précision@k, le rappel@k, le rang du premier passage pertinent et la fidélité de
la réponse finale.

La troisième comparaison concerne la génération de contenus. Une génération directe par
prompt unique peut être comparée à une génération structurée par plan et sections. Les
critères peuvent porter sur le respect du programme, la répétition, la longueur, la
cohérence de la progression, la qualité orale et le nombre de corrections nécessaires.

La quatrième comparaison concerne les reviews automatiques. Il faut mesurer ce qu'elles
détectent réellement : nombre de problèmes signalés, taux de corrections acceptées,
effets indésirables sur le texte et réduction de la charge de relecture humaine. Cette
évaluation est importante pour éviter de présenter les reviews comme une garantie de
qualité sans preuve.

La cinquième comparaison concerne l'impact opérationnel. Il faut comparer le temps de
production et le nombre d'interventions humaines avant et après automatisation. Même si
les chiffres restent partiels, cette mesure est indispensable pour relier la solution à la
problématique économique de l'entreprise.
